% 冷却与衰变——指数衰减模型
% 活动六：牛顿冷却（茶杯降温）；活动七：放射性半衰期与碳-14 测年。

\section{活动六：茶杯降温——冷却中的指数衰减}

\subsection{情境}
一杯 $90$℃ 的热水放在室温 $20$℃ 的房间里，测得 $8$ min 后水温为 $70$℃。

\subsection{提出问题}
照这样下去，水大约什么时候降到 $40$℃？如果手头没有温度计，只凭几次水温读数，能把室温也反推出来吗？

\subsection{选模型}
牛顿冷却定律：水的散热速率与水和环境的温差成正比，
\[
  \frac{dT}{dt}=-k\,(T-T_{\text{env}})\quad\Longrightarrow\quad
  T(t)=T_{\text{env}}+(T_0-T_{\text{env}})e^{-kt}.
\]
温差 $\Delta T=T-T_{\text{env}}$ 按 $\Delta T(t)=\Delta T_0e^{-kt}$ 衰减，
所以每经过一个“半衰期” $\tau=\dfrac{\ln2}{k}$，温差就减半。
为什么不用一次函数？因为降温是先快后慢：散热快慢取决于当前的温差，
而不是匀速地每分钟降相同的度数。

\subsection{拟合}
由 $T_{\text{env}}=20$℃、$T(0)=90$℃、$T(8)=70$℃：
\[
  e^{-8k}=\frac{70-20}{90-20}=\frac{5}{7}
  \quad\Longrightarrow\quad
  k=\frac{\ln(7/5)}{8}\approx0.0421\ \mathrm{min}^{-1},
\]
温差半衰期
\[
  \tau=\frac{\ln2}{k}\approx16.5\ \mathrm{min}.
\]
实验记录如下（每 $4$ min 读一次温度）：

\begin{center}
\footnotesize
\setlength{\tabcolsep}{4pt}
\renewcommand{\arraystretch}{1.4}
\begin{tabular}{p{2.8cm}p{3.4cm}p{3.6cm}}
\toprule
时间 $t$（min） & 水温 $T$（℃） & 温差 $T-20$（℃） \\
\midrule
$0$ & $90.0$ & $70.0$ \\
$4$ & $79.2$ & $59.2$ \\
$8$ & $70.0$ & $50.0$ \\
$12$ & $62.3$ & $42.3$ \\
$16$ & $55.7$ & $35.7$ \\
$20$ & $50.2$ & $30.2$ \\
$24$ & $45.5$ & $25.5$ \\
$28$ & $41.6$ & $21.6$ \\
\bottomrule
\end{tabular}
\end{center}

预测降到 $40$℃ 的时刻：由 $e^{-kt}=\dfrac{40-20}{90-20}=\dfrac{2}{7}$，
\[
  t=\frac{\ln(7/2)}{k}\approx29.8\ \text{min},
\]
即大约 $30$ min 时降到 $40$℃ 以下（$30$ min 时约 $39.8$℃）。

\subsection{检验与反思}
用未参与拟合的时刻校验：模型给出 $T(20)\approx50.2$℃、$T(24)\approx45.5$℃、
$T(28)\approx41.6$℃，与表中读数一致，差异在四舍五入的 $0.1$℃ 之内。
表中温差列也吻合半衰期结构：温差从 $70$℃ 减半到 $35$℃ 只需约
$16.5$ min，$16$ min 时实测 $35.7$℃ 正好落在附近。

若连室温也不知道，可用等间隔的三次读数 $T_1,T_2,T_3$ 反推：
由 $(T_2-T_{\text{env}})^2=(T_1-T_{\text{env}})(T_3-T_{\text{env}})$ 解得
\[
  T_{\text{env}}=\frac{T_2^2-T_1T_3}{2T_2-T_1-T_3},
\]
代入 $90.0$、$70.0$、$55.7$ 得约 $19.8$℃，与温度计读数一致。

模型的边界：温差指数衰减意味着水温可以无限接近 $20$℃ 但不会更低，
把曲线外推到 $0$℃ 以下没有物理意义；实际过程还受蒸发、杯壁吸热、
空气流动等因素影响。
\key{指数衰减的是“差值”：温差每过半衰期减半，半衰期是冷却过程的“指纹”。}

\section{活动七：碳-14 测年——半衰期与年代}

\subsection{情境}
活着的生物不断与大气交换碳，体内碳-14 的比例保持稳定；死亡之后交换停止，
碳-14 以半衰期 $5730$ 年持续衰变。考古队采回古木和骨片，想知道它们距今多少年。

\subsection{提出问题}
测出样品中碳-14 的残留比例 $r=\dfrac{N}{N_0}$，怎样换算成年代？
这个方法能覆盖多古老的样品？

\subsection{选模型}
放射性衰变的快慢与现存原子核数成正比：$\dfrac{dN}{dt}=-\lambda N$，于是
\[
  N(t)=N_0\left(\frac{1}{2}\right)^{t/T},\qquad T=5730\ \text{年},
\]
衰变常数 $\lambda=\dfrac{\ln2}{T}\approx1.21\times10^{-4}$（每年）。
由残留比例反推年代：
\[
  t=\frac{T}{\ln2}\ln\frac{N_0}{N}=\frac{5730}{\ln2}\ln\frac{1}{r}
  \approx8267\ln\frac{1}{r}.
\]
为什么不用线性模型？因为衰变是“按比例”进行的：每个原子核在某段时间内
衰变的概率大致固定，剩下的越少，单位时间衰变掉的也越少。

\subsection{拟合}
盖革计数器记下 $10$ min 内的总次数，减去本底即为样品贡献的净计数。
设现代碳每克每分钟约 $15$ 次，本底 $10$ min 内 $5$ 次，则
$r=\dfrac{\text{净计数率}/\text{质量}}{15}$。四组测量如下：

\begin{center}
\footnotesize
\setlength{\tabcolsep}{4pt}
\renewcommand{\arraystretch}{1.4}
\begin{tabular}{p{2.2cm}p{1.5cm}p{2.0cm}p{1.8cm}p{2.4cm}p{1.8cm}p{2.2cm}}
\toprule
样品 & 质量（g） & 总计数（10 min） & 本底（10 min） & 净计数率（次/(min·g)） & 残留比例 $r$ & 年代（年） \\
\midrule
现代碳标准 & $2.0$ & $305$ & $5$ & $15.0$ & $1.00$ & $0$ \\
木炭甲 & $2.0$ & $155$ & $5$ & $7.5$ & $0.50$ & $5730$ \\
骨片乙 & $2.0$ & $101$ & $5$ & $4.8$ & $0.32$ & 约 $9400$ \\
古墓木炭丙 & $4.0$ & $80$ & $5$ & $1.875$ & $0.125$ & $17190$ \\
\bottomrule
\end{tabular}
\end{center}

例如骨片乙：$r=0.32$，
\[
  t=\frac{5730}{\ln2}\ln\frac{1}{0.32}\approx9419\ \text{年}\approx9400\ \text{年}.
\]
木炭甲的 $r=0.50$ 恰为一个半衰期（$5730$ 年），古墓木炭丙的
$r=0.125=1/8$ 恰为三个半衰期（$17190$ 年），与公式互相印证。

\subsection{检验与反思}
适用区间：样品越老，残留比例越小。$5$ 万年时
$r=0.5^{50000/5730}\approx0.24\%$，信号已淹没在本底涨落里；
太年轻的样品 $r$ 接近 $1$，又难与污染区分。
所以碳-14 测年一般适用于几百年到 $5$ 万年。

两个主要风险：
\begin{itemize}
  \item 污染：混入现代碳（防腐材料、手汗、空气中的二氧化碳）使 $r$ 偏大、
  年代偏轻；混入古老碳（煤尘、石灰岩）使 $r$ 偏小、年代偏重。
  例如距今约 $1.7$ 万年的样品若混入 $1\%$ 的现代碳，$r$ 约从 $0.125$
  升到 $0.134$，测算年代会从 $17190$ 年变成约 $16600$ 年，年轻了约 $600$ 年。
  \item 本底计数：老样品净计数很小，本底的微小偏差会被放大。
  例如古墓木炭丙的本底若多计 $1$ 次，净计数从 $75$ 降为 $74$，$r$ 从
  $0.125$ 变为 $0.123$，年代从 $17190$ 年变为约 $17300$ 年。
\end{itemize}
所以报出“年代”的同时，必须给出误差范围和样品处理记录。
\key{半衰期把“残留比例”翻译成“年代”，但比例太小或受到污染时，翻译就会失真。}
